\documentclass[preprint]{aastex631}
\usepackage{hyperref}

\begin{document}

\title{Research Notes}
\author{Agastya Gaur}

\section{PS Dataset}
The datset was produced by the NASA Exoplanet Archive \href{
http://exoplanetarchive.ipac.caltech.edu}{(link)}.

In cases of duplicate rows per host star, the row flagged as default parameters was retained. Information about default parameters can be found \href{https://exoplanetarchive.ipac.caltech.edu/docs/faq.html}{here}. Information about the stellar data cols can be found \href{https://exoplanetarchive.ipac.caltech.edu/docs/API_STELLARHOSTS_columns.html}{here}.

A composite parameters table was NOT used, as composite parameters are not guaranteed to be self-consistent. PCA on stellar parameters assumes that each data vector represents a coherent physical model of the star. Mixing references can cause inconsistencies.

For multi-star systems, only the star listed as the planet host was retained.

\section{Data Imputation}
The data matrix $X$ contains a significant amount of NaN values.
\begin{verbatim}
0/4506 NaN in col 0
536/4506 NaN in col 1
629/4506 NaN in col 2
341/4506 NaN in col 3
1228/4506 NaN in col 4
4408/4506 NaN in col 5
3444/4506 NaN in col 6
779/4506 NaN in col 7
2111/4506 NaN in col 8
3175/4506 NaN in col 9
3553/4506 NaN in col 10
4121/4506 NaN in col 11
3979/4506 NaN in col 12
\end{verbatim}

A Multiple Imputation by Chained Equations (MICE) approach is being considered for handling this. Possible references: \citep{campbellFullyBayesianApproach2024}

\bibliography{references}

\end{document}